\endinput
%------------------------------------------------------------------------------%
\begin{question}
\end{question}

%------------------------------------------------------------------------------%
\begin{resolution}{Solução}
\end{resolution}

%------------------------------------------------------------------------------%
\endinput

\begin{frame}{Exemplo -- Laplace em uma matriz $4\times4$}
Calcule
\[
\det(A),
\qquad
A=
\begin{pmatrix}
1&0&0&2\\
0&3&0&0\\
4&0&5&0\\
0&0&0&2
\end{pmatrix}.
\]
\medskip

A primeira coluna possui dois zeros, mas a quarta coluna possui
três zeros. Expandimos pela quarta coluna:
\[
\det(A)
=
2C_{44}.
\]
\medskip

Como $4+4$ é par,
\[
C_{44}=M_{44}.
\]
\medskip

Logo,
\[
\det(A)
=
2
\begin{vmatrix}
1&0&0\\
0&3&0\\
4&0&5
\end{vmatrix}.
\]
\end{frame}

\begin{frame}{Exemplo -- Laplace em uma matriz $4\times4$}
A matriz restante é triangular inferior:
\[
\begin{pmatrix}
1&0&0\\
0&3&0\\
4&0&5
\end{pmatrix}.
\]
Portanto,
\[
M_{44}=1\cdot3\cdot5=15.
\]
Assim,
\[
\det(A)=2\cdot15=30.
\]
\medskip

Logo,
\[
\boxed{\det(A)=30}.
\]
\end{frame}
